\subsection{Playing Codenames}

% Prompt Engineering ChatGPT for Codenames

Sidji and Stephenson evaluate how prompt engineering affects GPT-4 performance in Codenames, a cooperative word game where a Codemaster gives one-word clues to help a Guesser identify team words while avoiding a hidden assassin word. They test five Codemaster prompt strategies across 50 games each, with the Guesser always using the Default strategy. The five strategies are: Default (a straightforward prompt), Cautious (avoid risky clues), Risky (be aggressive), Chain-of-Thought (reason step by step before answering), and Self-Refine (generate an answer, critique it, then produce a final answer).

The paper was easy to follow. Strategies were clearly described with full prompt text, and both the implementation and the underlying Codenames framework were publicly available. The main challenge was cost---replicating experiments with GPT-4 was prohibitive, so we used GPT-4o-mini and Gemini 2.5 Flash Lite instead. We also noticed the Guesser strategy is fixed throughout the original study, which left open whether Guesser prompting affects outcomes at all.

We extended the design by testing all 36 Codemaster $\times$ Guesser strategy combinations across 30 fixed board seeds (1,101 OpenAI games, 664 Gemini games). Guesser strategy turned out to matter far more than Codemaster strategy: varying the Guesser produced a 53-percentage-point spread in win rate (29\% -- 82\%), compared to only 10 points when varying the Codemaster. The best combination, Chain-of-Thought Codemaster paired with a Cautious Guesser, hit 96\% win rate, exceeding the original paper's best result of 94\% on a cheaper model. We also found a systematic failure in the Self-Refine Guesser: a critique prompt asking whether the guess ``could accidentally be the assassin'' caused the model to select the assassin in 85\% of games. This happened because the Guesser has no access to hidden role information---when asked to reason about which word might be the assassin, it would guess, convince itself a particular word was dangerous, and then switch to it thinking it was being careful. Reframing the critique positively reduced assassin hits to 66\% and tripled the win rate. More broadly, our Codemaster results reproduce the original finding that prompt strategy has limited effect, but show that this conclusion is an artifact of holding the Guesser fixed---once the Guesser is allowed to vary, strategy choice clearly matters.
